%-----------------------------------------------------------------------------------
%   PACKAGES 
%-----------------------------------------------------------------------------------

\documentclass[a4paper,11pt,abstracton]{scrartcl}

\usepackage[margin=.9in]{geometry}


\usepackage{mlmodern}


\usepackage{amsmath,amsthm,amssymb}
\usepackage{mathtools,mathrsfs}
\usepackage{centernot}
\usepackage{wrapfig}

\usepackage{tikz,tikz-cd,asymptote}
\usetikzlibrary{positioning, fit, calc}

\usepackage[most]{tcolorbox}
\usepackage{hyperref,wrapfig}
\usepackage[dvipsnames]{xcolor}
\usepackage{enumerate}
\usepackage[a]{esvect}
\usepackage{tocloft}


\usepackage{sectsty}
\makeatletter
\def\@seccntformat#1{\textcolor{Red}{\hspace*{0.2em}$\textbf{\S}$\hspace{-0.3em} \csname the#1\endcsname}\hspace{0.5em}}
\makeatother

\usepackage{fancyhdr}
\pagestyle{fancy}
\lhead{Analytical Number Theory}
\rhead{\thepage}
\cfoot{} 

\setlength{\parindent}{0cm}

%-----------------------------------------------------------------------------------
%   COMMANDS
%-----------------------------------------------------------------------------------

\DeclarePairedDelimiter{\floor}{\lfloor}{\rfloor}

\newcommand{\nonadj}{\,\neg\textsf{adj}\,}
\newcommand{\adj}{\,\textsf{adj}\,}
\newcommand{\RR}{\mathbb{R}}
\newcommand{\ZZ}{\mathbb{Z}}
\newcommand{\NN}{\mathbb{N}}
\newcommand{\CC}{\mathbb{C}}
\newcommand{\QQ}{\mathbb{Q}}
\newcommand{\sgn}{\text{sgn\hspace{0.5mm}}}
\newcommand{\osc}{\text{osc\hspace{0.5mm}}}
\newcommand{\foralls}{\hspace{2mm}\forall}
\newcommand{\rimply}{\Rightarrow}
\newcommand{\limply}{\Leftarrow}
\newcommand{\underdash}{\rule{0.25cm}{0.1mm}\hspace{0.5mm}}
\newcommand{\undersmalldash}{\rule{0.18cm}{0.1mm}\hspace{0.3mm}}
\newcommand{\s}{\square}
\newcommand{\xr}{\xrightarrow}

\newcommand{\genstirlingI}[3]{%
  \genfrac{[}{]}{0pt}{#1}{#2}{#3}%
}
\newcommand{\genstirlingII}[3]{%
  \genfrac{\{}{\}}{0pt}{#1}{#2}{#3}%
}
\newcommand{\genlah}[3]{%
  \genfrac{\lfloor}{\rfloor}{0pt}{#1}{#2}{#3}%
}
\newcommand{\stirlingI}[2]{\genstirlingI{}{#1}{#2}}
\newcommand{\stirlingII}[2]{\genstirlingII{}{#1}{#2}}
\newcommand{\lah}[2]{\genlah{}{#1}{#2}}

%-----------------------------------------------------------------------------------
%   ENVIROMENTS
%-----------------------------------------------------------------------------------

\newtheoremstyle{dotlessP}{}{}{}{}{\color{RedViolet}\bfseries}{.}{ }{}
\theoremstyle{dotlessP}
\newtheorem{problem}{Problem}
\newtheoremstyle{dotlessS}{}{}{}{}{\color{ForestGreen}\bfseries}{.}{ }{}
\theoremstyle{dotlessS}
\newtheorem*{solution}{Solution}

\theoremstyle{remark}
\newtheorem*{remark}{Remark}
\theoremstyle{definition}
\newtheorem*{theorem}{Theorem}

\definecolor{brandblue}{rgb}{0.34, 0.7, 1}

\newtcolorbox{mainbox}[1]{
  enhanced jigsaw,
  breakable,
  arc=0mm,
  bottomtitle=0.5mm,
  boxrule=0mm,
  colbacktitle=black!10!white,
  colframe=brandblue,
  coltitle=black,
  fonttitle=\bfseries,
  left=2.5mm,
  leftrule=1mm,
  right=3.5mm,
  title={#1},
  toptitle=0.75mm
}

\newtcolorbox{subbox}[1]{
  arc=0mm,
  bottomtitle=0.5mm,
  boxrule=0mm,
  colbacktitle=black!10!white,
  colframe=black!30!white,
  coltitle=black,
  fonttitle=\bfseries,
  left=2.5mm,
  leftrule=1mm,
  right=3.5mm,
  title={#1},
  toptitle=0.75mm
}

\newtcolorbox{definition_box}{
enhanced jigsaw,
breakable,
pad at break*=1mm,
colback=lightgray!12!white,
colframe=RedViolet!45!black,
colbacktitle=Gray!38!white,
before={\vspace{4mm}},
after={\vspace{2mm}},
sharp corners
}


\newtcolorbox{problem_box}{
enhanced jigsaw,
breakable,
pad at break*=1mm,
colback=white!5!white,
colframe=gray!75!black,
colbacktitle=white!28!white,
sharp corners,
before={\vspace{6mm}},
after={\vspace{2mm}}
}

\newtcolorbox{problem_box_1}{
enhanced jigsaw,
breakable,
pad at break*=1mm,
colback=SeaGreen!5!white,
colframe=SeaGreen!75!black,
colbacktitle=Black!28!white,
sharp corners,
before={\vspace{6mm}},
after={\vspace{2mm}}
}

\newtcolorbox{problem_box_title}[1]{
enhanced jigsaw,
breakable,
pad at break*=1mm,
colback=Apricot!5!white,
colframe=Melon!75!black,
sharp corners,
before={\vspace{2mm}},
after={\vspace{2mm}},
title={#1}, 
fonttitle=\bfseries
}

\newtcolorbox{claim_box}{
enhanced jigsaw,
breakable,
pad at break*=1mm,
colback=white,
colframe=white,
borderline west={2pt}{0pt}{white},
sharp corners,
before={\vspace{1mm}},
after={\vspace{0.8mm}}
}

\newtcolorbox{example_box}{
enhanced jigsaw,
breakable,
pad at break*=1mm,
colback=lightgray!12!white,
colframe=lightgray!15!white,
borderline west={2pt}{0pt}{black},
sharp corners,
before={\vspace{2mm}},
after={\vspace{2mm}}
}

\newtcolorbox{remark_box}{
enhanced jigsaw,
breakable,
pad at break*=1mm,
colback=Apricot!5!white,
colframe=Melon!75!black,
sharp corners,
before={\vspace{2mm}},
after={\vspace{2mm}}
}


\begin{document}

%-----------------------------------------------------------------------------------
%   TOP MATTERS
%-----------------------------------------------------------------------------------

\title{Analytical Number Theory}
\subtitle{Handkerchief}
\author{Pratyay Sen}
\date{\today}
\maketitle

\tableofcontents

\eject


%-----------------------------------------------------------------------------------
%   MAIN BODY
%-----------------------------------------------------------------------------------

\section{Some Basic Notation}
We begin with some basic notations that we will be using throughout this document. These notations are essential to analytical number theory in the study of number theoretic functions.
Let $f,g$ be two functions. Let $c, k$ be any constants. Then we say,
\begin{enumerate}
\item $f=O(g) \iff |f|<cg$
\item $f=o(g) \iff f/g \rightarrow 0$
\item $f \sim g \iff f/g \rightarrow 1$
\item $f \prec g \iff f/g \rightarrow 0 \iff f=o(g)$
\item $f \succ g \iff f/g \rightarrow \infty$
\item $f \asymp g \iff cg < f < kg$
\end{enumerate}

\begin{mainbox}{Growth estimates for common arithmetic functions:}
For every fixed $\varepsilon>0$ and every fixed $k\geq0$, the
following estimates hold as $n\to\infty$:
\[
    \mathbf{1}(n)=O(1),
    \qquad
    \mu(n)=O(1),
    \qquad
    \lambda(n)=O(1),
\]
\[
    d(n)=O(n^{\varepsilon}),
    \qquad
    \varphi(n)=O(n),
    \qquad
    \operatorname{id}(n)=O(n),
\]
\[
    \sigma_k(n)=O(n^{k+\varepsilon}),
    \qquad
    \omega(n)=O(\log n),
    \qquad
    \Omega(n)=O(\log n),
\]
and
\[
    \Lambda(n)=O(\log n)
\]
\end{mainbox}

\begin{claim_box}
\textit{Proof that $d(n)=O(n^{\varepsilon})$:}
Fix $\varepsilon>0$, and write
\[
    n=\prod_{p\mid n}p^{\alpha_p}
\]
Then
\[
    d(n)=\prod_{p\mid n}(\alpha_p+1)
\]
For every integer $\alpha\geq1$,
\[
    \alpha+1\leq2^{\alpha}
\]
Thus, whenever $p\geq2^{1/\varepsilon}$,
\[
    \alpha+1
      \leq2^{\alpha}
      \leq p^{\varepsilon\alpha}
\]

Only finitely many primes satisfy $p<2^{1/\varepsilon}$. For each of
these primes, let
\[
    C_{p,\varepsilon}
      =\max_{\alpha\geq0}
        \frac{\alpha+1}{p^{\varepsilon\alpha}}
\]
This maximum is finite because
$(\alpha+1)p^{-\varepsilon\alpha}\to0$ as $\alpha\to\infty$.
Therefore
\begin{align*}
    \frac{d(n)}{n^{\varepsilon}}
      &=\prod_{p\mid n}
        \frac{\alpha_p+1}{p^{\varepsilon\alpha_p}}\\
      &\leq\prod_{p<2^{1/\varepsilon}}C_{p,\varepsilon}
       =:C_{\varepsilon}
\end{align*}
The constant $C_{\varepsilon}$ is finite and independent of $n$.
Hence
\[
    d(n)\leq C_{\varepsilon}n^{\varepsilon},
\]
which proves the result.
\hfill$\square$
\end{claim_box}

\section{Euler--Maclaurin Summation and the Sum-to-Integral Method}

Many sums in analytic number theory can be estimated by replacing the
sum with an integral and controlling the resulting error. The
Euler--Maclaurin summation formula makes this procedure precise. Its
first-order form is often already sufficient for obtaining the main
term of a summatory function.

\subsection{The Sawtooth Function}

\begin{definition_box}
For $t\in\mathbb{R}$, let
\[
    \{t\}=t-\floor{t}
\]
be the fractional part of $t$, and define the \emph{sawtooth function}
by
\[
    \psi(t)=\{t\}-\frac12
\]
The function $\psi$ is periodic with period $1$ and satisfies
\[
    -\frac12\leq\psi(t)<\frac12
\]
Moreover, for every integer $m$,
\[
    \int_m^{m+1}\psi(t)\,dt=0
\]
\end{definition_box}

\subsection{Euler's Summation Formula}
\label{subsec:euler-summation}

The motivation behind this formula that looks absurd at first, is the we need a way to be able to approximate sums in some way. Now see that sums have an innate connection with integrals. Consider a function $f(x)$ and plot its graph, then if you want to visualise the sum, it would be the area of several rectangles whose y-component is $f(n)$ and 1 is the x-component. Thus the sum is the area under several rectangles of width 1 unit and height $f(n)$. This area is kind of close to the integral of $f(x)$, which leads us to this formula.

\begin{mainbox}{First-order Euler--Maclaurin formula:}
Let $A<B$ be integers, and let $f$ be continuously differentiable on
$[A,B]$. Then
\[
    \sum_{n=A}^{B}f(n)
      =\int_A^B f(t)\,dt
       +\frac{f(A)+f(B)}2
       +\int_A^B\psi(t)f'(t)\,dt
\]
\end{mainbox}

\begin{claim_box}
\textit{Proof:}
On each interval $[m,m+1]$, we have
\[
    \psi(t)=t-m-\frac12
\]
Integration by parts gives
\begin{align*}
    \int_m^{m+1}\psi(t)f'(t)\,dt
      &=\left[\left(t-m-\frac12\right)f(t)\right]_m^{m+1}
        -\int_m^{m+1}f(t)\,dt\\
      &=\frac{f(m)+f(m+1)}2
        -\int_m^{m+1}f(t)\,dt
\end{align*}
Summing this identity over the integers $m$ satisfying
$A\leq m\leq B-1$ gives
\[
    \int_A^B\psi(t)f'(t)\,dt
      =\frac{f(A)+f(B)}2
       +\sum_{n=A+1}^{B-1}f(n)
       -\int_A^Bf(t)\,dt
\]
Rearranging proves the formula.
\hfill$\square$
\end{claim_box}

\begin{mainbox}{Sum-to-integral estimate:}
Under the same assumptions,
\[
    \sum_{n=A}^{B}f(n)
      =\int_A^B f(t)\,dt
       +\frac{f(A)+f(B)}2
       +O\!\left(\int_A^B|f'(t)|\,dt\right)
\]
If $f$ is monotone on $[A,B]$, then
\[
    \sum_{n=A}^{B}f(n)
      =\int_A^B f(t)\,dt
       +O\bigl(|f(A)|+|f(B)|\bigr)
\]
In particular, if $f$ is positive and decreasing, the error is
$O(f(A))$, while if $f$ is positive and increasing, the error is
$O(f(B))$.
\end{mainbox}

\begin{claim_box}
\textit{Proof:}
Since $|\psi(t)|\leq1/2$, the first-order formula gives
\[
    \left|\int_A^B\psi(t)f'(t)\,dt\right|
      \leq\frac12\int_A^B|f'(t)|\,dt
\]
If $f$ is monotone, then
\[
    \int_A^B|f'(t)|\,dt=|f(B)-f(A)|,
\]
and the stated bounds follow.
\hfill$\square$
\end{claim_box}

\begin{example_box}
Taking $f(t)=\log t$ gives
\begin{align*}
    \sum_{n\leq x}\log n
      &=\int_1^x\log t\,dt+O(\log x)\\
      &=x\log x-x+O(\log x)
\end{align*}
More generally, for fixed $\alpha>-1$,
\[
    \sum_{n\leq x}n^{\alpha}
      =\frac{x^{\alpha+1}}{\alpha+1}
       +
       \begin{cases}
           O(x^{\alpha}) & \text{if $\alpha\geq0$},\\
           O(1) & \text{if $-1<\alpha<0$}
       \end{cases}
\]
Replacing $x$ by $\floor{x}$ changes only the stated error terms.
\end{example_box}


\subsection{The Harmonic Sum}
\label{subsec:harmonic-sum}

\begin{mainbox}{Asymptotic formula for the harmonic sum:}
For $x\geq2$,
\begin{equation}
\label{eq:harmonic-sum}
    \sum_{n\leq x}\frac1n
      =\log x+\gamma+O\!\left(\frac1x\right),
\end{equation}
where $\gamma$ is the Euler--Mascheroni constant. In particular,
\[
    \sum_{n\leq x}\frac1n=\log x+O(1)
\]
\end{mainbox}

\begin{claim_box}
\textit{Proof:}
Let $N$ be a positive integer and apply the first-order formula to
$f(t)=1/t$ on $[1,N]$. Since $f'(t)=-1/t^2$, we obtain
\[
    \sum_{n=1}^{N}\frac1n
      =\log N+\frac12\left(1+\frac1N\right)
       -\int_1^N\frac{\psi(t)}{t^2}\,dt
\]
The improper integral converges absolutely because
$|\psi(t)|\leq1/2$. Define
\[
    \gamma
      =\frac12-\int_1^{\infty}\frac{\psi(t)}{t^2}\,dt
\]
It follows that
\[
    \sum_{n=1}^{N}\frac1n-\log N-\gamma
      =\frac1{2N}
       +\int_N^{\infty}\frac{\psi(t)}{t^2}\,dt
\]
Therefore
\begin{align*}
    \left|
      \sum_{n=1}^{N}\frac1n-\log N-\gamma
    \right|
      &\leq\frac1{2N}
        +\frac12\int_N^{\infty}\frac{dt}{t^2}\\
      &\leq\frac1N
\end{align*}
Thus
\[
    \sum_{n=1}^{N}\frac1n
      =\log N+\gamma+O\!\left(\frac1N\right)
\]
Finally, taking $N=\floor{x}$ gives
\[
    \log N=\log x+O\!\left(\frac1x\right),
\]
which proves the result for real $x\geq2$.
\hfill$\square$
\end{claim_box}

\begin{remark_box}
A simpler integral comparison already gives the weaker estimate needed
in many arguments:
\[
    \log(N+1)
      \leq\sum_{n=1}^{N}\frac1n
      \leq1+\log N
\]
Hence the harmonic sum differs from $\log N$ by a bounded quantity.
Euler--Maclaurin identifies the limiting constant and improves the
error to $O(1/N)$.
\end{remark_box}

\section{Arithmetic Functions}
\label{sec:arithmetic-functions}
We have proved several results regarding arithmetic functions in the "Elementary Number Theory" document. We shall list the functions below:
\\For
\[
    n=\prod_{i=1}^{r}p_i^{\alpha_i},
\]
the $p_i$ are distinct primes and the $\alpha_i$ are positive integers.
\begin{enumerate}
\item \textbf{The constant-one function:}
\[
    \mathbf{1}(n)=1
\]

\item \textbf{The identity function:}
\[
    \operatorname{id}(n)=n
\]

\item \textbf{The divisor-counting function:}
\[
    d(n)=\tau(n)=\sum_{d\mid n}1
      =\prod_{i=1}^{r}(\alpha_i+1)
\]

\item \textbf{The sum-of-divisors functions:}
\[
    \sigma_k(n)=\sum_{d\mid n}d^k
\]
The ordinary sum-of-divisors function is
\[
    \sigma(n)=\sigma_1(n)=\sum_{d\mid n}d
      =\prod_{i=1}^{r}
        \frac{p_i^{\alpha_i+1}-1}{p_i-1}
\]

\item \textbf{Euler's totient function:}
\[
    \varphi(n)=\#\{1\leq a\leq n:\gcd(a,n)=1\}
      =n\prod_{p\mid n}\left(1-\frac1p\right)
\]

\item \textbf{The Möbius function:}
\[
    \mu(n)=
      \begin{cases}
          1 & \text{if $n=1$},\\
          (-1)^r & \text{if $n=\prod_{i=1}^{r}p_i$ is squarefree},\\
          0 & \text{if a prime square divides $n$}
      \end{cases}
\]
Its absolute value $|\mu(n)|$ is the indicator function of the
squarefree integers.

\item \textbf{The Dirichlet delta function:}
\[
    \delta(n)=
      \begin{cases}
          1 & \text{if $n=1$},\\
          0 & \text{if $n>1$}
      \end{cases}
\]
It is the identity element for Dirichlet convolution.

\item \textbf{The distinct-prime-factor function:}
\[
    \omega(n)=r,
\]
the number of distinct primes dividing $n$, with $\omega(1)=0$.

\item \textbf{The total-prime-factor function:}
\[
    \Omega(n)=\sum_{i=1}^{r}\alpha_i,
\]
the number of prime factors of $n$ counted with multiplicity, with
$\Omega(1)=0$.

\item \textbf{The von Mangoldt function:}
\[
    \Lambda(n)=
      \begin{cases}
          \log p & \text{if $n=p^m$ for a prime $p$ and $m\geq1$},\\
          0 & \text{otherwise}
      \end{cases}
\]

\item \textbf{Liouville's function:}
\[
    \lambda(n)=(-1)^{\Omega(n)}
\]
Thus $\lambda$ is completely multiplicative and $\lambda(1)=1$.

\item \textbf{The logarithm as an arithmetic function:}
It appears naturally in the convolution identity
\[
    \Lambda*\mathbf{1}= \text{log}(n)
\]
\end{enumerate}

\subsection{Dirichlet Convolution and Dirichlet Series}
\label{subsec:dirichlet-convolution}

\begin{definition_box}
Let $f,g:\mathbb{N}\to\mathbb{C}$ be arithmetic functions. Their
\emph{Dirichlet convolution} is the arithmetic function $f*g$ defined by
\[
    (f*g)(n)=\sum_{d\mid n}f(d)g\!\left(\frac nd\right)
\]
Equivalently,
\[
    (f*g)(n)=\sum_{ab=n}f(a)g(b),
\]
where the second sum is over all ordered pairs of positive integers
whose product is $n$.
\end{definition_box}

\begin{example_box}
For example,
\[
    (f*g)(6)
      =f(1)g(6)+f(2)g(3)+f(3)g(2)+f(6)g(1)
\]
Some important convolution identities are
\[
    \mu*\mathbf{1}=\delta,
    \qquad
    \mathbf{1}*\mathbf{1}=d,
    \qquad
    \operatorname{id}*\mathbf{1}=\sigma,
    \qquad
    \varphi*\mathbf{1}=\operatorname{id}
\]
\end{example_box}

\begin{mainbox}{Möbius formula for Euler's totient function:}
For every positive integer $n$,
\begin{equation}
\label{eq:totient-mobius}
    \varphi(n)
      =(\operatorname{id}*\mu)(n)
      =n\sum_{d\mid n}\frac{\mu(d)}d
\end{equation}
\end{mainbox}

\begin{claim_box}
\textit{Proof:}
Using
\[
    \varphi*\mathbf{1}=\operatorname{id},
    \qquad
    \mathbf{1}*\mu=\delta,
\]
associativity of Dirichlet convolution gives
\begin{align*}
    \varphi
      &=\varphi*\delta
       =\varphi*(\mathbf{1}*\mu)\\
      &=(\varphi*\mathbf{1})*\mu
       =\operatorname{id}*\mu
\end{align*}
Consequently,
\begin{align*}
    \varphi(n)
      &=\sum_{d\mid n}d\,\mu\!\left(\frac nd\right)\\
      &=\sum_{d\mid n}\frac nd\,\mu(d)
       =n\sum_{d\mid n}\frac{\mu(d)}d
\end{align*}
\hfill$\square$
\end{claim_box}

\begin{definition_box}
Let $f$ be an arithmetic function. Its \emph{Dirichlet series} is
\[
    D_f(s)=\sum_{n=1}^{\infty}\frac{f(n)}{n^s},
    \qquad s=\sigma+it\in\mathbb{C},
\]
whenever the series converges. Here
\[
    n^{-s}=e^{-s\log n}
\]
The series converges absolutely at $s$ when
\[
    \sum_{n=1}^{\infty}\frac{|f(n)|}{n^\sigma}<\infty
\]
\end{definition_box}

\begin{remark_box}
Let $s=\sigma+it$. If
\[
    \sum_{n=1}^{\infty}\frac{|f(n)|}{n^\sigma}<\infty,
\]
then the Dirichlet series
\[
    D_f(s)=\sum_{n=1}^{\infty}\frac{f(n)}{n^s}
\]
converges absolutely, and therefore converges.
\end{remark_box}

\begin{claim_box}
\textit{Proof:}
Since $s=\sigma+it$, we have
\[
    |n^{-s}|=|e^{-s\log n}|
      =e^{-\sigma\log n}=n^{-\sigma}
\]
Consequently,
\[
    \left|\frac{f(n)}{n^s}\right|
      =\frac{|f(n)|}{n^\sigma}
\]
Thus the assumed convergence is exactly the absolute convergence of
$D_f(s)$. More explicitly, for $M>N$,
\[
    \left|\sum_{n=N+1}^{M}\frac{f(n)}{n^s}\right|
      \leq\sum_{n=N+1}^{M}\frac{|f(n)|}{n^\sigma}
\]
The right-hand side tends to $0$ as $N,M\to\infty$, because the
series of absolute values converges. Hence the partial sums of
$D_f(s)$ form a Cauchy sequence in $\mathbb{C}$, so $D_f(s)$ converges.
\hfill$\square$
\end{claim_box}

\begin{remark_box}
If $f(n)=O(n^A)$ for some real number $A$, then $D_f(s)$ converges
absolutely in the half-plane
\[
    \operatorname{Re}(s)>A+1,
\]
because its absolute values are bounded by a constant multiple of the
convergent series $\sum n^{A-\sigma}$.
\end{remark_box}

\begin{mainbox}{Convolution becomes multiplication:}
Suppose the Dirichlet series of $f$ and $g$ converge absolutely at
$s$. Then
\[
    D_{f*g}(s)=D_f(s)D_g(s)
\]
\end{mainbox}

\begin{claim_box}
\textit{Proof:}
Absolute convergence allows us to multiply and rearrange the two
series:
\begin{align*}
    D_f(s)D_g(s)
      &=\left(\sum_{a=1}^{\infty}\frac{f(a)}{a^s}\right)
        \left(\sum_{b=1}^{\infty}\frac{g(b)}{b^s}\right)\\
      &=\sum_{a,b\geq1}\frac{f(a)g(b)}{(ab)^s}\\
      &=\sum_{n=1}^{\infty}\frac{1}{n^s}
        \sum_{ab=n}f(a)g(b)\\
      &=\sum_{n=1}^{\infty}\frac{(f*g)(n)}{n^s}
       =D_{f*g}(s)
\end{align*}
\hfill$\square$
\end{claim_box}

\begin{mainbox}{Divisor sums and Dirichlet series:}
Suppose
\[
    f(n)=\sum_{d\mid n}g(d)
\]
At every $s$ for which $D_g(s)$ and $\zeta(s)$ converge absolutely,
\[
    D_f(s)=D_g(s)\zeta(s)
\]
\end{mainbox}

\begin{claim_box}
\textit{Proof:}
Since the constant-one function satisfies $\mathbf{1}(n)=1$, we have
\[
    f(n)=\sum_{d\mid n}g(d)\mathbf{1}\!\left(\frac nd\right)
      =(g*\mathbf{1})(n)
\]
The convolution-product formula therefore gives
\[
    D_f(s)
      =D_{g*\mathbf{1}}(s)
      =D_g(s)D_{\mathbf{1}}(s)
      =D_g(s)\zeta(s)
\]
Equivalently, absolute convergence allows the direct rearrangement
\begin{align*}
    D_f(s)
      &=\sum_{n=1}^{\infty}\frac1{n^s}
        \sum_{d\mid n}g(d)\\
      &=\sum_{d,m\geq1}\frac{g(d)}{(dm)^s}\\
      &=\left(\sum_{d=1}^{\infty}\frac{g(d)}{d^s}\right)
        \left(\sum_{m=1}^{\infty}\frac1{m^s}\right)
       =D_g(s)\zeta(s)
\end{align*}
\hfill$\square$
\end{claim_box}

\subsection{Generalised Convolutions}

Dirichlet convolution combines two arithmetic functions at a fixed
integer. In analytic number theory, it is also useful to convolve an
arithmetic function with a function of a real variable.

\begin{definition_box}
Let $\alpha:\mathbb{N}\to\mathbb{C}$ be an arithmetic function and
let $f:[1,\infty)\to\mathbb{C}$. Their \emph{generalised convolution}
is the function
\[
    (\alpha\circ f)(x)
      :=\sum_{n\leq x}\alpha(n)f\!\left(\frac xn\right),
    \qquad x\geq1
\]
The sum is over all positive integers $n\leq x$. Unlike a Dirichlet
series, this is a finite sum for every fixed $x$.
\end{definition_box}

\begin{example_box}
If $f(x)=1$, then
\[
    (\alpha\circ f)(x)=\sum_{n\leq x}\alpha(n),
\]
which is the summatory function of $\alpha$. If $f(x)=x^s$, then
\[
    (\alpha\circ f)(x)
      =x^s\sum_{n\leq x}\frac{\alpha(n)}{n^s},
\]
so generalised convolutions also naturally produce truncated
Dirichlet series.
\end{example_box}

\begin{remark_box}
Generalised convolution is linear in both variables. For constants
$a,b\in\mathbb{C}$,
\[
    ((a\alpha+b\beta)\circ f)(x)
      =a(\alpha\circ f)(x)+b(\beta\circ f)(x),
\]
and
\[
    (\alpha\circ(af+bg))(x)
      =a(\alpha\circ f)(x)+b(\alpha\circ g)(x)
\]
\end{remark_box}

\begin{mainbox}{Associativity with Dirichlet convolution:}
Let $\alpha$ and $\beta$ be arithmetic functions, and let
$f:[1,\infty)\to\mathbb{C}$. Then, for every $x\geq1$,
\[
    \alpha\circ(\beta\circ f)
      =(\alpha*\beta)\circ f
\]
\end{mainbox}

\begin{claim_box}
\textit{Proof:}
By the definition of generalised convolution,
\begin{align*}
    (\alpha\circ(\beta\circ f))(x)
      &=\sum_{n\leq x}\alpha(n)
        (\beta\circ f)\!\left(\frac xn\right)\\
      &=\sum_{n\leq x}\alpha(n)
        \sum_{m\leq x/n}\beta(m)
        f\!\left(\frac{x}{nm}\right)\\
      &=\sum_{nm\leq x}\alpha(n)\beta(m)
        f\!\left(\frac{x}{nm}\right)
\end{align*}
All these sums are finite, so we may group the pairs $(n,m)$ according
to their product $r=nm$. This gives
\begin{align*}
    (\alpha\circ(\beta\circ f))(x)
      &=\sum_{r\leq x}
        \left(\sum_{nm=r}\alpha(n)\beta(m)\right)
        f\!\left(\frac xr\right)\\
      &=\sum_{r\leq x}(\alpha*\beta)(r)
        f\!\left(\frac xr\right)\\
      &=((\alpha*\beta)\circ f)(x)
\end{align*}
Therefore
\[
    \alpha\circ(\beta\circ f)=(\alpha*\beta)\circ f
\]
\hfill$\square$
\end{claim_box}

\begin{mainbox}{Generalised inversion formula:}
Let $\alpha$ be an arithmetic function with $\alpha(1)\neq0$, so that
its Dirichlet inverse $\alpha^{-1}$ exists. For functions
$F,G:[1,\infty)\to\mathbb{C}$,
\[
    G=\alpha\circ F
      \quad\Longleftrightarrow\quad
      F=\alpha^{-1}\circ G
\]
\end{mainbox}

\begin{claim_box}
\textit{Proof:}
Since $\alpha^{-1}$ is the Dirichlet inverse of $\alpha$,
\[
    \alpha^{-1}*\alpha
      =\alpha*\alpha^{-1}=\delta
\]
Also, the delta function is the identity for generalised convolution,
because
\[
    (\delta\circ F)(x)
      =\sum_{n\leq x}\delta(n)F\!\left(\frac xn\right)
      =F(x)
\]

Suppose first that $G=\alpha\circ F$. Using the associativity formula,
\[
    \alpha^{-1}\circ G
      =\alpha^{-1}\circ(\alpha\circ F)
      =(\alpha^{-1}*\alpha)\circ F
      =\delta\circ F=F
\]
Thus $F=\alpha^{-1}\circ G$.

Conversely, suppose that $F=\alpha^{-1}\circ G$. Then
\[
    \alpha\circ F
      =\alpha\circ(\alpha^{-1}\circ G)
      =(\alpha*\alpha^{-1})\circ G
      =\delta\circ G=G
\]
Hence $G=\alpha\circ F$, which proves both directions.
\\Note that we can also substitute $\alpha^{-1}(n)$ by $\mu(n)\alpha(n)$.
\hfill$\square$
\end{claim_box}

\begin{mainbox}{Sum with Dirichlet convolution:}
Let $\alpha$ and $\beta$ be arithmetic functions, and define the
summatory function
\[
    B(x)=\sum_{m\leq x}\beta(m)
\]
Then
\[
    (\alpha\circ B)(x)
      =\sum_{r\leq x}(\alpha*\beta)(r)
\]
\end{mainbox}

\begin{claim_box}
\textit{Proof:}
Expanding the definition of $B$ gives
\begin{align*}
    (\alpha\circ B)(x)
      &=\sum_{n\leq x}\alpha(n)
        B\!\left(\frac xn\right)\\
      &=\sum_{n\leq x}\alpha(n)
        \sum_{m\leq x/n}\beta(m)\\
      &=\sum_{nm\leq x}\alpha(n)\beta(m)
\end{align*}
Grouping the pairs $(n,m)$ according to their product $r=nm$, we
obtain
\[
    \sum_{nm\leq x}\alpha(n)\beta(m)
      =\sum_{r\leq x}\sum_{nm=r}\alpha(n)\beta(m)
      =\sum_{r\leq x}(\alpha*\beta)(r)
\]
\hfill$\square$
\end{claim_box}

\begin{example_box}
Taking $\beta=\mathbf{1}$ gives $B(x)=\floor{x}$, and therefore
\[
    \sum_{n\leq x}\alpha(n)\floor{\frac xn}
      =\sum_{r\leq x}(\alpha*\mathbf{1})(r)
\]
\end{example_box}

\subsection{The Riemann Zeta Function}
\label{subsec:riemann-zeta}

\begin{definition_box}
The \emph{Riemann zeta function} is initially defined in the half-plane
$\operatorname{Re}(s)>1$ by
\[
    \zeta(s)=\sum_{n=1}^{\infty}\frac1{n^s}
\]
It is the Dirichlet series of the constant-one function:
\[
    \zeta(s)=D_{\mathbf{1}}(s)
\]
\end{definition_box}

\begin{mainbox}{Euler product:}
For $\operatorname{Re}(s)>1$,
\[
    \zeta(s)=\prod_{p}
      \left(1-\frac1{p^s}\right)^{-1},
\]
where the product is over all prime numbers.
\end{mainbox}

\begin{claim_box}
\textit{Proof:}
For each prime $p$ and $\operatorname{Re}(s)>1$, the geometric series
gives
\[
    \left(1-\frac1{p^s}\right)^{-1}
      =\sum_{j=0}^{\infty}\frac1{p^{js}}
\]
Multiplying these series over a finite set of primes produces one term
$n^{-s}$ for every positive integer $n$ whose prime divisors all lie
in that set. By unique prime factorization, each such integer occurs
exactly once. Letting the finite set increase to include every prime,
and using the absolute convergence of $\sum n^{-s}$ when
$\operatorname{Re}(s)>1$, we obtain
\[
    \prod_p\left(1-\frac1{p^s}\right)^{-1}
      =\sum_{n=1}^{\infty}\frac1{n^s}
      =\zeta(s)
\]
\hfill$\square$
\end{claim_box}

\begin{mainbox}{Euler's evaluation of $\zeta(2)$:}
\begin{equation}
\label{eq:zeta-two}
    \zeta(2)=\sum_{n=1}^{\infty}\frac1{n^2}
      =\frac{\pi^2}{6}
\end{equation}
We use this without proof.
\end{mainbox}


\begin{example_box}
For $\operatorname{Re}(s)>1$,
\[
    \sum_{n=1}^{\infty}\frac{d(n)}{n^s}=\zeta(s)^2,
    \qquad
    \sum_{n=1}^{\infty}\frac{\mu(n)}{n^s}=\frac1{\zeta(s)},
\]
while for $\operatorname{Re}(s)>2$,
\[
    \sum_{n=1}^{\infty}\frac{\sigma(n)}{n^s}
      =\zeta(s)\zeta(s-1),
    \qquad
    \sum_{n=1}^{\infty}\frac{\varphi(n)}{n^s}
      =\frac{\zeta(s-1)}{\zeta(s)}
\]
\end{example_box}

\begin{claim_box}
\textit{Proof:}
Recall the convolution identities
\[
    d=\mathbf{1}*\mathbf{1},
    \qquad
    \mu*\mathbf{1}=\delta,
    \qquad
    \sigma=\operatorname{id}*\mathbf{1},
    \qquad
    \varphi*\mathbf{1}=\operatorname{id}
\]
Also,
\[
    D_{\mathbf{1}}(s)=\zeta(s),
    \qquad
    D_\delta(s)=1,
\]
and
\[
    D_{\operatorname{id}}(s)
      =\sum_{n=1}^{\infty}\frac{n}{n^s}
      =\sum_{n=1}^{\infty}\frac1{n^{s-1}}
      =\zeta(s-1)
\]

For $\operatorname{Re}(s)>1$, the Dirichlet series of
$\mathbf{1}$ and $\mu$ converge absolutely, since
$|\mu(n)|\leq1$. Therefore the convolution-product formula gives
\[
    \sum_{n=1}^{\infty}\frac{d(n)}{n^s}
      =D_d(s)
      =D_{\mathbf{1}*\mathbf{1}}(s)
      =D_{\mathbf{1}}(s)^2
      =\zeta(s)^2
\]
Similarly,
\[
    D_\mu(s)\zeta(s)
      =D_{\mu*\mathbf{1}}(s)
      =D_\delta(s)=1,
\]
so
\begin{equation}
\label{eq:mu-dirichlet-series}
    \sum_{n=1}^{\infty}\frac{\mu(n)}{n^s}
      =D_\mu(s)=\frac1{\zeta(s)}
\end{equation}
In particular, this also shows that $\zeta(s)\neq0$ when
$\operatorname{Re}(s)>1$.

For $\operatorname{Re}(s)>2$, the series of
$\operatorname{id}$, $\sigma$, and $\varphi$ converge absolutely:
indeed, $\varphi(n)\leq n$ and
\[
    \sigma(n)=\sum_{d\mid n}d
      \leq n\sum_{d\mid n}\frac1d
      \leq n\sum_{d=1}^{n}\frac1d
      \leq n(1+\log n)
\]
Hence
\[
    \sum_{n=1}^{\infty}\frac{\sigma(n)}{n^s}
      =D_{\operatorname{id}*\mathbf{1}}(s)
      =D_{\operatorname{id}}(s)D_{\mathbf{1}}(s)
      =\zeta(s-1)\zeta(s),
\]
and
\[
    D_\varphi(s)\zeta(s)
      =D_{\varphi*\mathbf{1}}(s)
      =D_{\operatorname{id}}(s)
      =\zeta(s-1)
\]
Since $\zeta(s)\neq0$ in this region, division gives
\[
    \sum_{n=1}^{\infty}\frac{\varphi(n)}{n^s}
      =D_\varphi(s)=\frac{\zeta(s-1)}{\zeta(s)}
\]
This proves all four identities.
\hfill$\square$
\end{claim_box}

\begin{mainbox}{Dirichlet series of the von Mangoldt function:}
For $\operatorname{Re}(s)>1$,
\[
    D_\Lambda(s)
      =\sum_{n=1}^{\infty}\frac{\Lambda(n)}{n^s}
      =-\frac{\zeta'(s)}{\zeta(s)}
\]
Thus the Dirichlet series of $\Lambda$ is the negative logarithmic
derivative of the Riemann zeta function.
\end{mainbox}

\begin{claim_box}
\textit{Proof:}
For every $n\geq1$, we have
\[
    0\leq\Lambda(n)\leq\log n
\]
Since $\sum (\log n)n^{-\sigma}$ converges for $\sigma>1$, the
Dirichlet series of $\Lambda$ converges absolutely in this half-plane.

On every closed half-plane
$\operatorname{Re}(s)\geq1+\varepsilon$, the series
\[
    \sum_{n=1}^{\infty}\frac{\log n}{n^s}
\]
converges uniformly by the Weierstrass $M$-test. We may therefore
differentiate the zeta series term by term:
\[
    \zeta'(s)
      =\sum_{n=1}^{\infty}\frac{d}{ds}(n^{-s})
      =-\sum_{n=1}^{\infty}\frac{\log n}{n^s}
\]
Hence
\[
    D_{\log}(s)
      =\sum_{n=1}^{\infty}\frac{\log n}{n^s}
      =-\zeta'(s)
\]

The convolution identity
\[
    \Lambda*\mathbf{1}=\log
\]
and the convolution-product formula now give
\[
    D_\Lambda(s)\zeta(s)
      =D_\Lambda(s)D_{\mathbf{1}}(s)
      =D_{\Lambda*\mathbf{1}}(s)
      =D_{\log}(s)
      =-\zeta'(s)
\]
Since $\zeta(s)\neq0$ for $\operatorname{Re}(s)>1$, division by
$\zeta(s)$ yields
\[
    D_\Lambda(s)
      =\sum_{n=1}^{\infty}\frac{\Lambda(n)}{n^s}
      =-\frac{\zeta'(s)}{\zeta(s)}
\]
\hfill$\square$
\end{claim_box}

\begin{mainbox}{Dirichlet series of Liouville's function:}
Let
\[
    \lambda(n)=(-1)^{\Omega(n)}
\]
For $\operatorname{Re}(s)>1$,
\[
    D_\lambda(s)
      =\sum_{n=1}^{\infty}\frac{\lambda(n)}{n^s}
      =\frac{\zeta(2s)}{\zeta(s)}
\]
\end{mainbox}

\begin{claim_box}
\textit{Proof:}
Since $|\lambda(n)|=1$, the Dirichlet series of $\lambda$ converges
absolutely when $\operatorname{Re}(s)>1$.

Let $q$ be the indicator function of the perfect squares:
\[
    q(n)=
      \begin{cases}
          1 & \text{if $n$ is a perfect square},\\
          0 & \text{otherwise}
      \end{cases}
\]
If
\[
    n=\prod_{i=1}^{r}p_i^{\alpha_i},
\]
then
\begin{align*}
    (\lambda*\mathbf{1})(n)
      &=\sum_{d\mid n}\lambda(d)\\
      &=\prod_{i=1}^{r}
        \left(\sum_{j=0}^{\alpha_i}(-1)^j\right)
\end{align*}
Each factor equals $1$ when $\alpha_i$ is even and $0$ when
$\alpha_i$ is odd. Hence
\[
    \lambda*\mathbf{1}=q
\]

The Dirichlet series of $q$ is
\[
    D_q(s)
      =\sum_{m=1}^{\infty}\frac1{(m^2)^s}
      =\zeta(2s)
\]
Applying the convolution-product formula gives
\[
    D_\lambda(s)\zeta(s)
      =D_{\lambda*\mathbf{1}}(s)
      =D_q(s)=\zeta(2s)
\]
Since $\zeta(s)\neq0$ for $\operatorname{Re}(s)>1$, we may divide by
$\zeta(s)$ and obtain
\[
    D_\lambda(s)=\frac{\zeta(2s)}{\zeta(s)}
\]
\hfill$\square$
\end{claim_box}

\begin{mainbox}{Dirichlet series of $k$-th-power-free integers:}
Fix an integer $k\geq2$, and define
\[
    q_k(n)=
      \begin{cases}
          1 & \text{if $n$ is $k$-th-power-free},\\
          0 & \text{otherwise}
      \end{cases}
\]
Thus $q_k(n)=1$ exactly when no prime power $p^k$ divides $n$.
For $\operatorname{Re}(s)>1$,
\[
    \sum_{n=1}^{\infty}\frac{q_k(n)}{n^s}
      =\frac{\zeta(s)}{\zeta(ks)}
\]
\end{mainbox}

\begin{claim_box}
\textit{Proof:}
Since $0\leq q_k(n)\leq1$, the Dirichlet series converges absolutely
for $\operatorname{Re}(s)>1$. The function $q_k$ is multiplicative,
and for every prime $p$,
\[
    q_k(p^j)=
      \begin{cases}
          1 & \text{if $0\leq j\leq k-1$},\\
          0 & \text{if $j\geq k$}
      \end{cases}
\]
Therefore its Euler factor at $p$ is the finite geometric sum
\[
    \sum_{j=0}^{\infty}\frac{q_k(p^j)}{p^{js}}
      =\sum_{j=0}^{k-1}\frac1{p^{js}}
      =\frac{1-p^{-ks}}{1-p^{-s}}
\]
Multiplying these factors over all primes gives
\begin{align*}
    \sum_{n=1}^{\infty}\frac{q_k(n)}{n^s}
      &=\prod_p\frac{1-p^{-ks}}{1-p^{-s}}\\
      &=\left(\prod_p(1-p^{-s})^{-1}\right)
        \left(\prod_p(1-p^{-ks})\right)\\
      &=\frac{\zeta(s)}{\zeta(ks)}
\end{align*}
\hfill$\square$
\end{claim_box}

\subsection{Logarithmic Derivatives}
\label{subsec:logarithmic-derivatives}

Let $F$ be a differentiable complex-valued function that does not
vanish in the region under consideration. Its \emph{logarithmic
derivative} is
\[
    \frac{F'(s)}{F(s)}
\]
Whenever a branch of $\log F(s)$ is defined, the chain rule gives
\[
    \frac{d}{ds}\log F(s)=\frac{F'(s)}{F(s)}
\]
The main advantage of the logarithmic derivative is that it converts
products into sums.

\begin{mainbox}{Product rule for logarithmic derivatives:}
If $F$ and $G$ are differentiable and nonzero, then
\[
    \frac{(FG)'(s)}{F(s)G(s)}
      =\frac{F'(s)}{F(s)}+\frac{G'(s)}{G(s)}
\]
Consequently, an absolutely and locally uniformly convergent product
$F(s)=\prod_jF_j(s)$ satisfies
\[
    \frac{F'(s)}{F(s)}
      =\sum_j\frac{F_j'(s)}{F_j(s)}
\]
whenever the differentiated sum converges locally uniformly.
\end{mainbox}

\begin{claim_box}
\textit{Proof:}
The ordinary product rule gives
\[
    (FG)'=F'G+FG'
\]
Dividing by $FG$ yields the first identity. The product formula follows
by applying the finite identity to partial products and passing to the
limit under the stated convergence assumptions.
\hfill$\square$
\end{claim_box}

\begin{remark_box}
Logarithmic derivatives record zeros and poles. If
$F(s)=(s-s_0)^mG(s)$ with $G(s_0)\neq0$, then
\[
    \frac{F'(s)}{F(s)}
      =\frac{m}{s-s_0}+\frac{G'(s)}{G(s)}
\]
Thus a zero of order $m$ gives a simple pole of the logarithmic
derivative with residue $m$. A pole of order $m$ gives residue $-m$.
\end{remark_box}

\begin{mainbox}{Arithmetic logarithmic differentiation:}
For an arithmetic function $f$, define
\[
    (\mathcal{D}f)(n)=f(n)\log n
\]
Then
\begin{equation}
\label{eq:arithmetic-log-derivative-product}
    \mathcal{D}(f*g)
      =(\mathcal{D}f)*g+f*(\mathcal{D}g)
\end{equation}
Whenever the Dirichlet series of $f$ may be differentiated term by
term,
\[
    D_{\mathcal{D}f}(s)=-D_f'(s)
\]
In particular,
\[
    \mathcal{D}\mathbf{1}=\log,
    \qquad
    \mathcal{D}\log=\log^2,
    \qquad
    \mathcal{D}\Lambda=\Lambda\log
\]
\end{mainbox}

\begin{claim_box}
\textit{Proof:}
For the convolution rule, use
$\log n=\log d+\log(n/d)$ to obtain
\begin{align*}
    \mathcal{D}(f*g)(n)
      &=\sum_{d\mid n}f(d)g\!\left(\frac nd\right)\log n\\
      &=\sum_{d\mid n}f(d)g\!\left(\frac nd\right)
        \left(\log d+\log\frac nd\right)\\
      &=((\mathcal{D}f)*g)(n)+(f*(\mathcal{D}g))(n)
\end{align*}
Also,
\[
    -D_f'(s)
      =-\frac{d}{ds}\sum_{n=1}^{\infty}\frac{f(n)}{n^s}
      =\sum_{n=1}^{\infty}\frac{f(n)\log n}{n^s}
      =D_{\mathcal{D}f}(s)
\]
The three special cases follow immediately from the definition.
\hfill$\square$
\end{claim_box}

\subsection{Selberg's Identity}
\label{subsec:selberg-identity}

\begin{mainbox}{Selberg's identity:}
For every positive integer $n$,
\begin{equation}
\label{eq:selberg-identity}
    \Lambda(n)\log n
      +\sum_{d\mid n}\Lambda(d)\Lambda\!\left(\frac nd\right)
      =\sum_{d\mid n}\mu(d)
        \left(\log\frac nd\right)^2
\end{equation}
Equivalently, in terms of Dirichlet convolution,
\[
    \Lambda\log+\Lambda*\Lambda=\mu*\log^2
\]
where $\Lambda\log$ denotes the pointwise product
$n\mapsto\Lambda(n)\log n$.
\end{mainbox}

\begin{claim_box}
\textit{Proof:}
We now use the fundamental convolution identity
\[
    \log=\Lambda*\mathbf{1}
\]
Applying $\mathcal{D}$ and using
Equation~\eqref{eq:arithmetic-log-derivative-product} gives
\begin{align*}
    \log^2
      &=\mathcal{D}(\log)\\
      &=\mathcal{D}(\Lambda*\mathbf{1})\\
      &=(\mathcal{D}\Lambda)*\mathbf{1}
        +\Lambda*(\mathcal{D}\mathbf{1})\\
      &=(\Lambda\log)*\mathbf{1}+\Lambda*\log
\end{align*}
Convolving both sides with $\mu$ and using associativity yields
\begin{align*}
    \mu*\log^2
      &=(\Lambda\log)*(\mu*\mathbf{1})
        +\Lambda*(\mu*\log)\\
      &=(\Lambda\log)*\delta+\Lambda*\Lambda\\
      &=\Lambda\log+\Lambda*\Lambda
\end{align*}
Here $\mu*\mathbf{1}=\delta$, while
$\mu*\log=\Lambda$ follows by Möbius inversion from
$\log=\Lambda*\mathbf{1}$. Evaluating the last convolution identity at
$n$ gives Equation~\eqref{eq:selberg-identity}.
\hfill$\square$
\end{claim_box}

\begin{remark_box}
Summing Equation~\eqref{eq:selberg-identity} over $n\leq x$ and
rearranging the finite divisor sums gives the exact summatory form
\begin{equation}
\label{eq:selberg-identity-summed}
    \sum_{n\leq x}\Lambda(n)\log n
      +\sum_{ab\leq x}\Lambda(a)\Lambda(b)
      =\sum_{dm\leq x}\mu(d)(\log m)^2
\end{equation}
This identity is the algebraic starting point for Selberg's symmetry
formula used in elementary proofs of the prime number theorem.
\end{remark_box}

\subsection{Ramanujan Sums}

Ramanujan sums are finite exponential sums taken over the reduced
residue classes modulo $n$. Although their definition involves complex
numbers, their values are always integers.

\begin{definition_box}
For a positive integer $n$ and an integer $k$, the \emph{Ramanujan sum} is
\[
    c_n(k)
      :=\sum_{\substack{1\leq a\leq n\\ \gcd(a,n)=1}}
        e^{2\pi iak/n}
\]
It depends only on the residue class of $k$ modulo $n$.
\end{definition_box}

\begin{example_box}
For $n=4$, the reduced residue classes are $1$ and $3$, so
\[
    c_4(k)=e^{2\pi ik/4}+e^{6\pi ik/4}
\]
In particular,
\[
    c_4(1)=0,
    \qquad
    c_4(2)=-2,
    \qquad
    c_4(4)=2
\]
\end{example_box}

\begin{mainbox}{Divisor formula for Ramanujan sums:}
For all positive integers $n$ and $k$,
\[
    c_n(k)=\sum_{d\mid\gcd(n,k)}
      d\,\mu\!\left(\frac nd\right)
\]
Consequently, $c_n(k)$ is an integer and depends only on
$\gcd(n,k)$.
\end{mainbox}

\begin{claim_box}
\textit{Proof:}
The Möbius divisor-sum identity gives
\[
    \sum_{d\mid m}\mu(d)
      =
      \begin{cases}
          1 & \text{if $m=1$},\\
          0 & \text{if $m>1$}
      \end{cases}
\]
Therefore the indicator of the condition $\gcd(a,n)=1$ is
\[
    \sum_{d\mid\gcd(a,n)}\mu(d)
\]
Substituting this into the definition of $c_n(k)$ and rearranging the
finite sums gives
\begin{align*}
    c_n(k)
      &=\sum_{a=1}^{n}e^{2\pi iak/n}
        \sum_{d\mid\gcd(a,n)}\mu(d)\\
      &=\sum_{d\mid n}\mu(d)
        \sum_{\substack{1\leq a\leq n\\ d\mid a}}
        e^{2\pi iak/n}
\end{align*}
Writing $a=db$, the inner sum becomes
\[
    \sum_{b=1}^{n/d}e^{2\pi ibk/(n/d)}
\]
This is the sum of all $(n/d)$-th roots of unity raised to the
$k$-th power, so it equals $n/d$ if $n/d\mid k$, and equals $0$
otherwise. Hence
\[
    c_n(k)
      =\sum_{\substack{d\mid n\\ n/d\mid k}}
        \mu(d)\frac nd
\]
Replacing $n/d$ by $q$ gives
\[
    c_n(k)
      =\sum_{\substack{q\mid n\\ q\mid k}}
        q\,\mu\!\left(\frac nq\right)
      =\sum_{q\mid\gcd(n,k)}
        q\,\mu\!\left(\frac nq\right)
\]
\hfill$\square$
\end{claim_box}

\begin{mainbox}{Closed formula for Ramanujan sums:}
Let $n$ be a positive integer, let $m$ be an integer, and put
\[
    \delta=\gcd(n,m)
\]
Then
\[
    c_n(m)=
      \mu\!\left(\frac n\delta\right)
      \frac{\varphi(n)}{\varphi(n/\delta)}
\]
Here $\delta$ denotes the integer $\gcd(n,m)$, not the Dirichlet delta
function.
\end{mainbox}

\begin{claim_box}
\textit{Proof:}
Write
\[
    n=\prod_{p\mid n}p^{a_p},
    \qquad
    \delta=\prod_{p\mid n}p^{b_p},
    \qquad 0\leq b_p\leq a_p
\]
Using the divisor formula and factoring the divisor sum prime by prime,
we obtain
\[
    c_n(m)
      =\prod_{p\mid n}
        \left(\sum_{j=0}^{b_p}
          p^j\mu\!\left(p^{a_p-j}\right)\right)
\]
For a fixed prime $p$, put $a=a_p$ and $b=b_p$. Since
$\mu(p^r)=0$ for $r\geq2$, the corresponding local factor is
\[
    \sum_{j=0}^{b}p^j\mu(p^{a-j})
      =
      \begin{cases}
          0 & \text{if $b\leq a-2$},\\
          -p^{a-1} & \text{if $b=a-1$},\\
          p^{a-1}(p-1) & \text{if $b=a$}
      \end{cases}
\]
These three values are exactly
\[
    \mu(p^{a-b})
      \frac{\varphi(p^a)}{\varphi(p^{a-b})},
\]
where $\varphi(1)=1$. Therefore
\begin{align*}
    c_n(m)
      &=\prod_{p\mid n}
        \mu(p^{a_p-b_p})
        \frac{\varphi(p^{a_p})}
             {\varphi(p^{a_p-b_p})}\\
      &=\mu\!\left(\prod_{p\mid n}p^{a_p-b_p}\right)
        \frac{\varphi(n)}
             {\varphi\!\left(\prod_{p\mid n}p^{a_p-b_p}\right)}\\
      &=\mu\!\left(\frac n\delta\right)
        \frac{\varphi(n)}{\varphi(n/\delta)}
\end{align*}
This proves the formula.
\hfill$\square$
\end{claim_box}

\begin{remark_box}
When $n\mid k$, and in particular when $k=n$, the divisor formula
becomes
\[
    c_n(k)=\sum_{d\mid n}d\,\mu\!\left(\frac nd\right)
      =(\operatorname{id}*\mu)(n)=\varphi(n)
\]
Thus the precise convolution identity is
\[
   c_n(n)=(\operatorname{id}*\mu)(n)=\varphi(n)
\]
For a general $k$, one must instead use the restricted divisor sum
over $d\mid\gcd(n,k)$.
\end{remark_box}

\begin{remark_box}
Two other useful special cases are
\[
    c_n(1)=\mu(n),
    \qquad
    c_n(0)=c_n(n)=\varphi(n)
\]
The first follows because $\gcd(n,1)=1$, while the second follows
because every exponential in the defining sum equals $1$.
\end{remark_box}

\section{Average Orders of Arithmetic Functions}
\label{sec:average-orders}

All the prerequisites for the following calculations have now been
proved. We shall use Euler summation from
Subsection~\ref{subsec:euler-summation}, the harmonic estimate
\eqref{eq:harmonic-sum}, the totient formula
\eqref{eq:totient-mobius}, the reciprocal zeta series
\eqref{eq:mu-dirichlet-series}, and the evaluation
\eqref{eq:zeta-two} of $\zeta(2)$.

\begin{definition_box}
Let $f$ and $g$ be arithmetic functions. We say that $g$ is an
\emph{average order} of $f$ if
\[
    \sum_{n\leq x}f(n)
      \sim\sum_{n\leq x}g(n)
      \qquad(x\to\infty)
\]
Thus an average order describes the typical size of $f(n)$ through
the asymptotic behaviour of its summatory function. Equivalently, 
$$\sum_{n\leq x}f(n)=xg(x)+o(xg(x))$$
\end{definition_box}

\begin{remark_box}
One must compare the two summatory functions. For example, if
$g(n)=cn$, then
\[
    \sum_{n\leq x}g(n)\sim\frac{c}{2}x^2,
\]
whereas for the slowly varying function $g(n)=\log n$,
\[
    \sum_{n\leq x}g(n)\sim xg(x)
\]
\end{remark_box}

\begin{mainbox}{Average orders of common arithmetic functions:}
Put
\[
    c=\zeta(2)=\frac{\pi^2}{6}
\]
by Equation~\eqref{eq:zeta-two}. Some standard average orders are
\begin{center}
\small
\renewcommand{\arraystretch}{1.35}
\setlength{\tabcolsep}{10pt}
\begin{tabular}{@{}lcc@{}}
    \hline
    {\textbf{Name}}
      &{\textbf{Function}}
      &{\textbf{Average order}}\\
    \hline
    Constant-one function
      & $\mathbf{1}(n)$ & $1$\\
    Identity function
      & $\operatorname{id}(n)$ & $n$\\
    Divisor-counting function
      & $d(n)$ & $\log n$\\
    Euler's totient function
      & $\varphi(n)$ & $\dfrac{n}{c}$\\
    Sum-of-divisors function
      & $\sigma(n)$ & $cn$\\
    Squarefree indicator
      & $|\mu(n)|$ & $\dfrac1c$\\
    $k$-th-power-free indicator, $k\geq2$
      & $q_k(n)$ & $\frac1{\zeta(k)}$\\
    Distinct-prime-factor function
      & $\omega(n)$ & $\log\log n$\\
    Total-prime-factor function
      & $\Omega(n)$ & $\log\log n$\\
    \hline
\end{tabular}
\end{center}
The average orders of $d$, $\varphi$, and $\sigma$ are proved below.
\end{mainbox}

\begin{mainbox}{Average order of the divisor function:}
The divisor-counting function $d(n)$ has average order $\log n$:
\[
    \sum_{n\leq x}d(n)\sim x\log x
\]
\end{mainbox}

\begin{claim_box}
\textit{Proof:}
Since $d(n)$ counts the factorizations $n=ab$ with $a,b\geq1$,
\[
    \sum_{n\leq x}d(n)
      =\sum_{a\leq x}\floor{\frac xa}
\]
Using $\floor{y}=y+O(1)$ and Equation~\eqref{eq:harmonic-sum},
\begin{align*}
    \sum_{n\leq x}d(n)
      &=x\sum_{a\leq x}\frac1a+O(x)\\
      &=x\log x+O(x)
\end{align*}
The logarithmic sum estimate proved in
Subsection~\ref{subsec:euler-summation} gives
\[
    \sum_{n\leq x}\log n
      =x\log x-x+O(\log x)
      \sim x\log x
\]
Hence
\[
    \sum_{n\leq x}d(n)
      \sim\sum_{n\leq x}\log n,
\]
so $\log n$ is an average order of $d(n)$.
\hfill$\square$
\end{claim_box}

\begin{mainbox}{Average order of Euler's totient function:}
Euler's totient function has average order
\[
    \frac{6}{\pi^2}n=\frac{n}{c}
\]
More precisely,
\[
    \sum_{n\leq x}\varphi(n)
      =\frac{3}{\pi^2}x^2+O(x\log x)
\]
\end{mainbox}

\begin{claim_box}
\textit{Proof:}
By Equation~\eqref{eq:totient-mobius},
\[
    \varphi(n)=n\sum_{d\mid n}\frac{\mu(d)}d
\]
Writing $n=dm$ and rearranging the finite sums gives
\[
    \sum_{n\leq x}\varphi(n)
      =\sum_{d\leq x}\mu(d)
        \sum_{m\leq x/d}m
\]
The power-sum estimate from
Subsection~\ref{subsec:euler-summation} gives
\[
    \sum_{m\leq y}m=\frac{y^2}{2}+O(y)
\]
Consequently, using Equation~\eqref{eq:harmonic-sum},
\[
    \sum_{n\leq x}\varphi(n)
      =\frac{x^2}{2}
        \sum_{d\leq x}\frac{\mu(d)}{d^2}
        +O(x\log x)
\]
Equation~\eqref{eq:mu-dirichlet-series} at $s=2$ and
Equation~\eqref{eq:zeta-two} give
\[
    \sum_{d=1}^{\infty}\frac{\mu(d)}{d^2}
      =\frac1{\zeta(2)}=\frac6{\pi^2}
\]
Moreover,
\[
    \sum_{d>x}\frac{|\mu(d)|}{d^2}
      \leq\sum_{d>x}\frac1{d^2}
      =O\!\left(\frac1x\right)
\]
Therefore
\[
    \sum_{n\leq x}\varphi(n)
      =\frac{x^2}{2\zeta(2)}+O(x\log x)
      =\frac3{\pi^2}x^2+O(x\log x)
\]
Since
\[
    \sum_{n\leq x}\frac6{\pi^2}n
      =\frac3{\pi^2}x^2+O(x),
\]
the function $6n/\pi^2=n/c$ is an average order of $\varphi(n)$.
\hfill$\square$
\end{claim_box}

\begin{mainbox}{Average order of the sum-of-divisors function:}
The sum-of-divisors function has average order
\[
    \frac{\pi^2}{6}n=cn
\]
More precisely,
\[
    \sum_{n\leq x}\sigma(n)
      =\frac{\pi^2}{12}x^2+O(x\log x)
\]
\end{mainbox}

\begin{claim_box}
\textit{Proof:}
Using the definition
\[
    \sigma(n)=\sum_{d\mid n}d
\]
from Section~\ref{sec:arithmetic-functions}, and writing $n=dm$, we
obtain
\[
    \sum_{n\leq x}\sigma(n)
      =\sum_{m\leq x}\sum_{d\leq x/m}d
\]
The power-sum estimate from
Subsection~\ref{subsec:euler-summation} gives
\[
    \sum_{d\leq y}d=\frac{y^2}{2}+O(y)
\]
Thus, using Equation~\eqref{eq:harmonic-sum},
\[
    \sum_{n\leq x}\sigma(n)
      =\frac{x^2}{2}\sum_{m\leq x}\frac1{m^2}
        +O(x\log x)
\]
By Equation~\eqref{eq:zeta-two} and
$\sum_{m>x}m^{-2}=O(1/x)$,
\begin{align*}
    \sum_{n\leq x}\sigma(n)
      &=\frac{\zeta(2)}2x^2+O(x\log x)\\
      &=\frac{\pi^2}{12}x^2+O(x\log x)
\end{align*}
Finally,
\[
    \sum_{n\leq x}\frac{\pi^2}{6}n
      =\frac{\pi^2}{12}x^2+O(x),
\]
so $\pi^2n/6=cn$ is an average order of $\sigma(n)$.
\hfill$\square$
\end{claim_box}

\section{Prime Number Function}

\begin{definition_box}
For every real number $x\geq0$, the \emph{prime-counting function} is
defined by
\[
    \pi(x)=\#\{p\leq x:p\text{ is prime}\}
\]
\end{definition_box}

Throughout this section, $\log$ denotes the natural logarithm.

\begin{mainbox}{An elementary lower bound:}
For every real number $x\geq2$,
\[
    \pi(x)>\log\log x
\]
\end{mainbox}

\begin{claim_box}
\textit{Proof:}
Let
\[
    p_1=2
\]
and let $p_n$ denote the $n$-th prime. We first prove by induction that
\[
    p_n\leq2^{2^{n-1}}\qquad(n\geq1)
\]
The case $n=1$ is the equality $p_1=2$. Suppose the estimate holds
for every $p_j$ with $1\leq j\leq n$. Every prime divisor of
\[
    N=\prod_{j=1}^{n}p_j+1
\]
is different from every $p_j$ with $1\leq j\leq n$. Therefore $N$ has a prime divisor
at least as large as $p_{n+1}$, and hence
\begin{align*}
    p_{n+1}
      &\leq \prod_{j=1}^{n}p_j+1\\
      &\leq 2^{\sum_{j=0}^{n-1}2^j}+1\\
      &=2^{2^n-1}+1
       \leq2^{2^n}
\end{align*}
This completes the induction.

Now let $n=\pi(x)$. Since $x\geq2$, we have $n\geq1$, and the next
prime satisfies $x<p_{n+1}$. Using the bound just proved,
\[
    x<p_{n+1}\leq2^{2^n}
\]
Taking logarithms twice gives
\[
    \log\log x
      <\log(2^n\log2)
      =n\log2+\log\log2
      <n,
\]
because $\log2<1$ and $\log\log2<0$. Since $n=\pi(x)$, this yields
\[
    \pi(x)>\log\log x
\]
\hfill$\square$
\end{claim_box}

\subsection{Prime Values of Polynomials}

\begin{mainbox}{No nonconstant polynomial is always prime:}
Let $f\in\mathbb{Z}[x]$ be a nonconstant polynomial. Then $f(n)$
cannot be prime for every positive integer $n$.
\end{mainbox}

\begin{claim_box}
\textit{Proof:}
Suppose, for a contradiction, that $f(n)$ is prime for every positive
integer $n$. Put
\[
    p=f(1)
\]
Then $p$ is prime. For every positive integer $k$, we have
\[
    1+kp\equiv1\pmod p
\]
Since a polynomial with integer coefficients preserves congruences,
\[
    f(1+kp)\equiv f(1)\equiv0\pmod p
\]
Thus $p\mid f(1+kp)$ for every $k\geq1$.

Because $f$ is nonconstant, $|f(x)|\to\infty$ as $x\to\infty$.
Choose $k$ sufficiently large that
\[
    f(1+kp)>p
\]
The integer $f(1+kp)$ is then divisible by the smaller positive integer
$p>1$, so it is composite. This contradicts the assumption that every
$f(n)$ is prime. Therefore no nonconstant polynomial in
$\mathbb{Z}[x]$ can be prime for every positive integer input.
\hfill$\square$
\end{claim_box}

\end{document}